%% =====================================================================
%%  B7-T-23-04 -- what B7 contributes to "The Gut Verdict"
%%  -------------------------------------------------------------------
%%  LAYER A: APPEND-ONLY. Written once at the entry's write, then left
%%  alone. If the source has more to say later, add a bullet -- do not
%%  rewrite. Material found ONLY here goes in the `unique` slot with
%%  the unique flag set, which makes it undeletable (rule R10).
%% =====================================================================
%% @kind: source
%% @id: B7-T-23-04
%% @book: B7
%% @topic: T-23-04
%% @locator: chs. 1--2 (pp. 13--22), additive
%% @status: distilled
%% @unique: no
%% @integrated: yes
%% @updated: 2026-09-19

\dossier{B7}{T-23-04}{\bookshort{B7}}{chs. 1--2 (pp. 13--22), additive}

\begin{distilled}
  \item B7 both confirms and disciplines the gut: distrust you cannot name is exactly where their readers start, but unexamined manner-reading is what they call global assessment --- guessing, not analysis.
  \item The social obstacle list explains why alarms get overridden: the belief that people will not lie to you, the training that calling someone a liar is worse than the act, the fear of the confrontation.
  \item The refinement: run the alarm forward into a question (the stimulus) and score the five seconds that follow --- the gut selects the moment; the window does the reading.
\end{distilled}
